\documentclass[a4paper]{article}

%% Language and font encodings
\usepackage[english]{babel}
\usepackage[utf8x]{inputenc}
\usepackage[T1]{fontenc}
\usepackage{ mathrsfs }
\usepackage{amsfonts}

%% Sets page size and margins
\usepackage[a4paper,top=2cm,bottom=2cm,left=3cm,right=3cm,marginparwidth=1.75cm]{geometry}

%% Useful packages
\usepackage{amsmath}
\usepackage{graphicx}
\usepackage[colorinlistoftodos]{todonotes}
\usepackage[colorlinks=true, allcolors=blue]{hyperref}
\title{Informal Notes on SVD of a matrix}
\author{Harsha Vardhan Tetali}
\begin{document}
\maketitle
\author

\section{Introduction}

In this article we consider finding the basis elements for the column space and row space of a matrix $A \in \mathbb{R}^{m \times n}$. Towards this objective, we first build positive-semi definite matrices which have row or column spaces equivalent to row and column spaces of $A$. We select positive semi-definite matrices for this task because positive semi-definite matrices are symmetric and can have an orthonormal eigen-basis which can span the required space.

\section{Examining the Fundamental Spaces}

\subsection{Null Spaces}
Let us start from the null space of matrix $A$.
Let $\boldsymbol{x} \in null(A)$.

Therefore,
\begin{equation}
A \boldsymbol{x} = \boldsymbol{0}_m
\end{equation}

where $\boldsymbol{0}_m \in \mathbb{R}^{m}$ is the zero vector in the $m$ dimensional space.

Multiplying $A^{\top}$ on both sides,

\begin{equation}
A^{\top} A \boldsymbol{x} = \boldsymbol{0}_n
\label{eq:ata}
\end{equation}

where $\boldsymbol{0}_n \in \mathbb{R}^{n}$ is the zero vector in the $n$ dimensional space.

Equation \ref{eq:ata} implies that $\boldsymbol{x} \in null(A^{\top}A)$.

Thus,

\begin{equation}
null(A) \subseteq null(A^{\top}A)
\label{eq:res1}
\end{equation}

Now let us say, $\boldsymbol{y} \in null(A^{\top}A)$.

Thus,

\begin{equation}
A^{\top}A \boldsymbol{y} = \boldsymbol{0}_n 
\end{equation}

Now pre-multiplying both sides by $\boldsymbol{y}^{\top}$ we have,

\begin{equation}
\boldsymbol{y}^{\top}A^{\top}A \boldsymbol{y} = 0
\end{equation}

This can be written as,

\begin{equation}
\left(A\boldsymbol{y}\right)^{\top}A \boldsymbol{y} = 0
\end{equation}

The above can be seen as a norm, therefore,

\begin{equation}
\| A \boldsymbol{y} \|_{2}^{2} = 0
\end{equation}

Thus, \begin{equation}
A \boldsymbol{y} = \boldsymbol{0}_m
\end{equation}

This implies, \begin{equation}
\boldsymbol{y} \subseteq null(A)
\end{equation}

Thus, 
\begin{equation}
null(A^{\top}A) \subseteq null(A)
\label{eq:res2}
\end{equation}

From \ref{eq:res1} and \ref{eq:res2} we have,

\begin{equation}
null(A^{\top}A) = null(A)
\end{equation}

Now playing the above tricks with $A^{\top}$ we have,

\begin{equation}
null(AA^{\top}) = null(A^{\top})
\label{eq:AAt}
\end{equation}

Thus, the null space of $A^{\top}A$ is the same as the null space of $A$ and the null space of $AA^{\top}$ is the same as null space of $A^{\top}$.

\subsection{Row and Column Spaces}

We know that the row space of $A$ is orthogonal to the null space of $A$ and from the relation between rank and nullity of a matrix, we can say that all spaces orthogonal to the row space is the null space. We therefore write this as,

\begin{equation}
R(A) = \left[null(A)\right]^{\perp}
\label{eq:usu}
\end{equation}

where $R(A)$ is the row-space of $A$.

Also,

\begin{equation}
R(A^{\top}) = \left[null(A^{\top})\right]^{\perp}
\end{equation}

This can be written as,

\begin{equation}
C(A) = \left[null(A^{\top})\right]^{\perp}
\end{equation}

where $C(A)$ is the column-space of $A$.

Combining the last equation with \ref{eq:AAt}, we have,

\begin{equation}
C(A) = \left[null(AA^{\top})\right]^{\perp}\label{eq:mine}
\end{equation}

But we know from \ref{eq:usu}, by putting $AA^{\top}$ instead of $A$, that,

\begin{equation}
R(AA^{\top}) = \left[null(AA^{\top})\right]^{\perp}
\end{equation}

Since, $R(AA^{\top})=C(AA^{\top})$, we can write the above as,

\begin{equation}
C(AA^{\top}) = \left[null(AA^{\top})\right]^{\perp}
\label{eq:yours}
\end{equation}

From \ref{eq:mine} and \ref{eq:yours} we have,

\begin{equation}
C(AA^{\top}) = C(A)
\end{equation}

Now replacing $A^{\top}$ instead of $A$, we have,

\begin{equation}
C(A^{\top}A) = C(A^{\top})
\end{equation}

As seen from the row space point of view, we have,

\begin{equation}
R(A^{\top}A) = R(A)
\end{equation}

Therefore, the column-space of $A$ is the same as the column-space of $AA^{\top}$. On similar lines, row-space of $A$ is same as the row-space of $A^{\top}A$.

\subsection{Basis for the Fundamental Spaces}

Let us look at the matrix $A \in \mathbb{R}^{m \times n}$ as a linear transformation of vectors from $\mathbb{R}^n$ to $\mathbb{R}^m$. Let us consider $m<n$. Thus we write,

$A: \mathbb{R}^n \rightarrow \mathbb{R}^m$

In this we have the domain set as the n-dimensional vector space $\mathbb{R}^n$ and the co-domain set as the m-dimensional vector space $\mathbb{R}^m$.

Similarly, the linear transformation, $A^{\top}: \mathbb{R}^m \rightarrow \mathbb{R}^n$ has the domain set as the m-dimensional vector space $\mathbb{R}^m$ and the co-domain set as the n-dimensional vector space $\mathbb{R}^n$.

Now observe that the spaces $C(A)$ and $null(A^{\top})$ form the entire space $\mathbb{R}^m$. Thus, $C(AA^{\top})$ and $null(AA^{\top})$ combined together form the entire space $\mathbb{R}^m$.

Since it is in the column space of $A$ it can be represented by the basis elements of the column space of $A$. This basis set can be easily found from the matrix $AA^{\top}$ since, it is a symmetric matrix we can find orthonormal basis by finding the eigenvectors of the symmetric matrix $AA^{\top}$. Let the orthonormal basis be $\{ \boldsymbol{u}_1, \boldsymbol{u}_2, \cdots, \boldsymbol{u}_r \}$. 

Now each $\boldsymbol{u}_i \in \mathbb{R}^m$, so let us see what happens to the basis under the transformation $A^{\top}$.

The new set, after transformation, is therefore, $\{ A^{\top}\boldsymbol{u}_1, A^{\top}\boldsymbol{u}_2, \cdots, A^{\top}\boldsymbol{u}_r \}$.

For any $i$, $j$ we have,

\begin{equation}
\left(A^{\top}\boldsymbol{u}_i \right)^{\top} A^{\top}\boldsymbol{u}_j = \boldsymbol{u}_i^{\top}AA^{\top}\boldsymbol{u}_j 
\label{eq:me}
\end{equation}


Since each is an eigenvector of $AA^{\top}$, we have

\begin{equation}
AA^{\top}\boldsymbol{u}_i = \lambda_i \boldsymbol{u}_i 
\label{eq:you}
\end{equation}

Therefore, from \ref{eq:me} and \ref{eq:you}, we have,

\begin{equation}
\boldsymbol{u}_i^{\top}AA^{\top}\boldsymbol{u}_j  = \boldsymbol{u}_i^{\top} \left( \lambda_j \boldsymbol{u}_j \right) = \lambda_j \boldsymbol{u}_i^{\top} \boldsymbol{u}_j = \begin{cases}
\lambda_i, \text{ where } i = j \\
0, \text{ where } i \neq j\\
\end{cases}
\label{eq:cas}
\end{equation}

Therefore, $\{ A^{\top}\boldsymbol{u}_1, A^{\top}\boldsymbol{u}_2, \cdots, A^{\top}\boldsymbol{u}_r \}$ is an orthogonal set.

Let $A^{\top}\boldsymbol{u}_i = \boldsymbol{w}_i$. From \ref{eq:cas} it is clear that, 
\begin{equation}
\boldsymbol{w}_i^{\top} \boldsymbol{w}_i = \lambda_i
\end{equation}

Let us define \begin{equation}
\boldsymbol{v}_i = \frac{\boldsymbol{w}_i}{\sqrt{\lambda_{i}} }
\end{equation} so as to normalize each $\boldsymbol{w}_i$ for all $i=1,2,\cdots,r$.

We therefore have,

\begin{equation}
A^{\top} \boldsymbol{u}_i = \sqrt{\lambda_{i}} \boldsymbol{v}_i
\end{equation}

Multiplying $A$ on both sides, we have,
\begin{equation}
AA^{\top} \boldsymbol{u}_i = \sqrt{\lambda_{i}} A\boldsymbol{v}_i
\end{equation}

From \ref{eq:you} and the above equation,

\begin{equation}
\lambda_i \boldsymbol{u}_i = \sqrt{\lambda_{i}} A\boldsymbol{v}_i
\end{equation}

Thus, we have,

\begin{equation}
A\boldsymbol{v}_i = \sqrt{\lambda_{i}}\boldsymbol{u}_i
\label{eq:sing}
\end{equation}

for all $i=1,2,\cdots,r$

We know that $\boldsymbol{v}_i \in \mathbb{R}^{n}$, therefore, we can add the remaining few vectors to the set $\{ \boldsymbol{v}_1, \boldsymbol{v}_2, \cdots, \boldsymbol{v}_r \}$, which is already an orthogonal set and therefore linearly independent, to create a basis for $\mathbb{R}^n$. Since, \begin{equation*}
\boldsymbol{v}_i = \frac{A^{\top}\boldsymbol{u}_i}{\sqrt{\lambda_i}}
\end{equation*}

We know that $\boldsymbol{v}_i$'s come from the column space of $A^{\top}$ and we also know that they form the basis of the columns space because they are linearly independent. Now we therefore consider orthonormal vectors that form the basis of null space of $A$. These can be obtained as the null space of $A^{\top}A$. These will be $n-r$ vectors, say, $\{ \boldsymbol{\eta}_1, \boldsymbol{\eta}_2, \cdots, \boldsymbol{\eta}_{n-r} \}$.

Therefore, \begin{equation}
A \boldsymbol{\eta}_j = \boldsymbol{0}
\label{eq:null}
\end{equation}

for all $j=1,2,\cdots,n-r$.

Putting vectors from \ref{eq:sing} and \ref{eq:null} as the column vectors of a matrix we have the following equality,

\begin{equation}
\begin{bmatrix}
A \boldsymbol{v}_1 & A \boldsymbol{v}_2 & \cdots & A \boldsymbol{v}_r & A \boldsymbol{\eta}_1 & A \boldsymbol{\eta}_2 & \cdots &A \boldsymbol{\eta}_{n-r}
\end{bmatrix} = \begin{bmatrix}
\sqrt{\lambda_{1}}\boldsymbol{u}_1 & \sqrt{\lambda_{2}}\boldsymbol{u}_2 &
\cdots \sqrt{\lambda_{r}}\boldsymbol{u}_r & \boldsymbol{0} & \boldsymbol{0} & \cdots & \boldsymbol{0}
\end{bmatrix}
\end{equation}

We can write the above as,

\begin{equation}
\begin{aligned}
A \begin{bmatrix}
\boldsymbol{v}_1 & \boldsymbol{v}_2 & \cdots & \boldsymbol{v}_r &  \boldsymbol{\eta}_1 & \boldsymbol{\eta}_2 & \cdots & \boldsymbol{\eta}_{n-r}
\end{bmatrix} = \\ \begin{bmatrix}
\boldsymbol{u}_1 & \boldsymbol{u}_2 & \cdots & \boldsymbol{u}_r
\end{bmatrix} \begin{bmatrix}
\sqrt{\lambda}_1 & 0 & \cdots &0 & 0 & 0 & \cdots & 0 \\
0 & \sqrt{\lambda}_2 & \cdots & 0 & 0 & 0 & \cdots & 0\\
\vdots & \vdots & \ddots & \vdots & \vdots & \vdots & \ddots & \vdots \\
0 & 0 & \cdots &\sqrt{\lambda}_r & 0 & 0 & \cdots & 0 \\
\end{bmatrix}
\end{aligned}
\end{equation}

From similar reasons, we can extend $\{
\boldsymbol{u}_1, \boldsymbol{u}_2, \cdots, \boldsymbol{u}_r
\}$ to form basis for $\mathbb{R}^m$. This can be done by adding the $m-r$ vectors that form the basis for the null space of $A^{\top}$. Let these be represented by $\{ \boldsymbol{\nu}_1, \boldsymbol{\nu}_2, \cdots, \boldsymbol{\nu}_{m-r} \}$. These can be obtained as the null space of $AA^{\top}$.

We can append these vectors in the matrix of the above equation:

\begin{equation}
\begin{aligned}
A \begin{bmatrix}
\boldsymbol{v}_1 & \boldsymbol{v}_2 & \cdots & \boldsymbol{v}_r &  \boldsymbol{\eta}_1 & \boldsymbol{\eta}_2 & \cdots & \boldsymbol{\eta}_{n-r}
\end{bmatrix} = \\
\begin{bmatrix}
\boldsymbol{u}_1 & \boldsymbol{u}_2 & \cdots & \boldsymbol{u}_r & \boldsymbol{\nu}_1 & \cdots & \boldsymbol{\nu}_{m-r}
\end{bmatrix} \begin{bmatrix}
\sqrt{\lambda_1} & 0 & \cdots &0 & 0 & 0 & \cdots & 0 \\
0 & \sqrt{\lambda_2} & \cdots & 0 & 0 & 0 & \cdots & 0\\
\vdots & \vdots & \ddots & \vdots & \vdots & \vdots & \ddots & \vdots \\
0 & 0 & \cdots & \sqrt{\lambda_r} & 0 & 0 & \cdots & 0 \\
0 & 0 & \cdots & 0 & 0 & 0 & \cdots & 0 \\
\vdots & \vdots & \ddots & \vdots & \vdots & \vdots & \ddots & \vdots \\
0 & 0 & \cdots & 0 & 0 & 0 & \cdots & 0 \\
\end{bmatrix}
\end{aligned}
\end{equation}

Let us denote,

\begin{equation}
V = \begin{bmatrix}
\boldsymbol{v}_1 & \boldsymbol{v}_2 & \cdots & \boldsymbol{v}_r &  \boldsymbol{\eta}_1 & \boldsymbol{\eta}_2 & \cdots & \boldsymbol{\eta}_{n-r}
\end{bmatrix} 
\end{equation}

And then denote,

\begin{equation}
U = \begin{bmatrix}
\boldsymbol{u}_1 & \boldsymbol{u}_2 & \cdots & \boldsymbol{u}_r & \boldsymbol{\nu}_1 & \cdots & \boldsymbol{\nu}_{m-r}
\end{bmatrix}
\end{equation}

Observe that, $\{ \boldsymbol{\eta}_1, \boldsymbol{\eta}_2, \cdots, \boldsymbol{\eta}_{n-r} \}$ can be chosen, using Gram-Schmidt Process, such that $V$ is an orthogonal matrix. Similarly, $U$ can also be chosen to be an orthogonal matrix. Therefore, we have,

\begin{equation}
U^{\top}U = \boldsymbol{I}_m
\end{equation}
and,
\begin{equation}
V^{\top}V = \boldsymbol{I}_n
\end{equation}

Let us also denote,

\begin{equation}
\Sigma = \begin{bmatrix}
\sqrt{\lambda_1} & 0 & \cdots &0 & 0 & 0 & \cdots & 0 \\
0 & \sqrt{\lambda_2} & \cdots & 0 & 0 & 0 & \cdots & 0\\
\vdots & \vdots & \ddots & \vdots & \vdots & \vdots & \ddots & \vdots \\
0 & 0 & \cdots & \sqrt{\lambda_r} & 0 & 0 & \cdots & 0 \\
0 & 0 & \cdots & 0 & 0 & 0 & \cdots & 0 \\
\vdots & \vdots & \ddots & \vdots & \vdots & \vdots & \ddots & \vdots \\
0 & 0 & \cdots & 0 & 0 & 0 & \cdots & 0 \\
\end{bmatrix}
\end{equation}

Therefore, we can write,

\begin{equation}
AV = U\Sigma 
\end{equation}

Thus,
\begin{equation}
A = U\Sigma V^{\top} 
\end{equation}

where $\Sigma \in \mathbb{R}^{m \times n}$ contains Singular Values of $A$, $U \in \mathbb{R}^{m \times m}$ and $V \in \mathbb{R}^{n \times n}$. 










\end{document}